\section{Introduction}

\begin{quote}
{\em Everyone who judges a paper cheaply reads only its title and
abstract.
Rite is a writing loop that makes the cheap read and the deep read
agree.}
\end{quote}

The point of research should not be the paper; it shold be the ideas
behind the papers.
Yet aloost uviersally, paper output is how we measure researcher
effectiveness: fields count their progress in paper generation and
ciations; careers, grants, and graduations are measured in terms of
paper poutput.
Which is strange since the paper is merely a receipt of effort spent,
not the goods.

And writing the receipt is slow.
The ceremony around a submission (searching, citing, formatting,
restating, self-checking) consumes months; an advisor who rewrites
every draft of every student loses those months many times over.
So an academic's life can run long on writing and short on the
research the writing is supposed to report.
Advice on that ceremony is plentiful and good: Shaw taxonomized the
questions reviewers ask of ICSE submissions~\cite{shaw03}; Widom gave
the five questions an introduction must answer~\cite{widom06};
Williams mapped working habits onto empirical
careers~\cite{williams11}; community standards now itemize what
reviewers may demand~\cite{ralph21}.
Yet all this advice shares one flaw: none of it is executable, so none
of it is checked, and applying it by hand is exactly the labor that
makes writing slow.

Until recently that flaw was incurable, since no tool could read a
literature, code it, and critique a draft.
Large language models changed that.
Newell's advice to a young scientist was to know your own strengths
and follow them~\cite{newell91}; a pipeline can now reverse-engineer
those strengths from a researcher's own recent papers, then map them
onto the problems a field currently ranks as most pressing.

So we built rite, which makes some of the old advice executable.
It consists of two small repositories: a read-only \textit{engine}
(github.com/timm/rite: a fourteen-step loop, eight prompt divisions,
seven Python scripts) and one \textit{workdir} per paper that holds
the goal, the coding vocabulary, and everything generated.
Per its own method, this paper is written by rite and paired with a
critique of itself.
This paper asks two questions.

\textbf{RQ1: Can the folklore of paper writing be encoded as a
checkable pipeline?}
Yes, partly.
Of the fourteen steps of Table~\ref{tab:loop}, we find \textbf{ten}
run as scripts or machine-checked gates; the irreducibly human residue
is small (steps 10 and 12, plus taste).

\textbf{RQ2: Is a cheap reading of the literature (title plus
abstract) a safe substitute for the full text?}
No, measurably.
In our exhibit, abstracts earned \textbf{9 of 21} full-text SE flags,
\textbf{3 of 10} multi-objective flags, and binary matching over full
text saturated on \textbf{21 of 22} papers until a frequency threshold
was applied (Section~\ref{sec:exhibit}).

Contributions: (1)~the loop and its eight prompt divisions, written so
a paper can quote the short form of each; (2)~a literature pipeline
whose byproducts (knee of the citation curve, download rate, coding
grids) become the paper's own gap argument; (3)~the
abstract-consistency gate: a submission's title-plus-abstract must
earn every topic flag its body earns; (4)~a replication package:
everything at github.com/timm/rite and github.com/timm/how2rite.

\subsection{FAQ}\label{sec:faq}

Three questions reach us before any others, so we answer them first.

\textbf{Doesn't encoding Newell-style problem
selection~\cite{newell76,newell91} steer a career into a rut?}
In our experience, no.
The literature step keeps forcing contact with material outside prior
experience, and the experiments that follow force a critical
evaluation of the tool base's premises.
The author offers his own last five years, coded by this same
pipeline, in the sequel to this draft; readers can judge the rut for
themselves.

\textbf{Will automation flood the field with average papers?}
Price's law observes that the square root of a field's contributors
produce half of its output~\cite{price63}.
The road into that square root is hard thinking, and no pipeline walks
it for you.
The tools here only cut the ceremony that surrounds the thinking.

\textbf{Does this kind of tooling ever find anything new?}
It did here.
While exploring this pipeline we found that a paper's title plus
abstract does not code the same way as its full text
(Section~\ref{sec:exhibit}).
That hazard has been on record since at least 2007, when Brereton et
al.\ warned that SE abstracts are too poor to rely on when selecting
primary studies~\cite{brereton07}.
But its consequences run far past review screening, since titles and
abstracts are what search engines index, what readers use to decide
what to read (so miscoding quietly corrupts manual literature
reviews), and what bidding systems consult when assigning reviewers to
papers.
Here, then, is a sizeable problem, barely discussed in SE, uncovered
by exactly the kind of LLM-assisted analysis that some researchers
would decline to run on methodological or ideological grounds.

\subsection{How to install}

\begin{verbatim}
git clone https://github.com/timm/rite
python3 rite/etc/init.py mypaper
cd mypaper
\end{verbatim}

The generated workdir README needs two lines filled in before anything
else runs; for example:

\begin{verbatim}
# mypaper
goal:  how to write SE research papers
years: 2021-2026
\end{verbatim}

Those two lines are the single source of truth for the whole
literature search: every downstream script reads its goal and year
range from here and nowhere else.

\subsection{How to use}

Follow Table~\ref{tab:loop}.
Steps 1--7 configure and search; step 8 gates entry (no elevator
speech, no paper); steps 9--14 loop until the self-test passes.
Each step names the prompt division that drives it, and each division
is a skill file whose short form fits on one screen.

\subsection{How to collaborate}

The repository is the paper; there is no other copy.
On every push to the main branch, continuous integration rebuilds the
PDF and serves it at a fixed URL (here,
timm.github.io/how2rite/paper1.pdf); pull requests only build, so
broken \TeX\ cannot land.
Merges stay small by construction: each section is one file under
\texttt{sec/}, and prose is one sentence per line, so two authors in
the same paragraph usually touch different lines.
Humans may commit to main; agents work on branches, submit pull
requests, edit only the section files an ownership table assigns them,
and never touch the preamble.
The built PDF is never committed; the served copy is the only copy.

\begin{table}
\caption{The rite loop. Divisions: ch=choose, se=search,
en=encode, st=structure, pr=prose, cr=critics, be=bench, sh=ship.
Steps marked $\ast$ are automatic.}
\label{tab:loop}
\small
\begin{tabular}{rlp{0.9cm}}
\toprule
\# & step & prompts\\
\midrule
1 & define your area of interest & ch\\
2 & define ``recent'' (field velocity) & se\\
3 & list your best recent papers & ch\\
4 & $\ast$ define what you are good at (from 3) & ch\\
5 & $\ast$ assemble critics (venues, standards,
    nearest ten) & cr\\
6 & $\ast$ search: knee, snowball; code full text
    at threshold & se,en\\
7 & define the most relevant problem (from 6),
    or the strangest & ch\\
8 & write paper[0]: title, abstract, elevator
    speech; else goto 7 & pr,st\\
9 & i=i+1; write paper[i]; rest stays headers & st,pr\\
10 & offline deep read, 30--60 minutes of
    comments & cr\\
11 & $\ast$ apply other critics; split fixes
    auto/manual & cr\\
12 & discuss with other humans & cr\\
13 & $\ast$ apply auto fixes; recode own abstract,
    must match body & en,pr\\
14 & $\ast$ self-test; fail goto 9, pass ship with
    replication package & sh\\
\bottomrule
\end{tabular}
\end{table}

\begin{figure*}
\centering
\begin{BVerbatim}[fontsize=\scriptsize]
 ENGINE github.com/timm/rite (read only)          WORKDIR one repo per paper
 +------------------------------------+           +---------------------------+
 | HOWTO.md   the 14 step loop        |  init.py  | README.md  goal:  years:  |
 | .claude/skills/  7 divisions       | --------> | flags.py   coding facets  |
 | etc/  fetch snowball stubs code    |           | bench.md   field norms    |
 |       getpdfs recode init          |  scripts  | TODO.md    work orders    |
 | pdf/  style source papers          | <-------- | lit/       all generated  |
 +------------------------------------+  run from | paper.tex  paper.pdf      |
                                          workdir  +---------------------------+
\end{BVerbatim}
\caption{The obligatory system diagram.
The argument it walks: advice lives once, on the left, read-only;
everything a paper generates lands on the right, so ten papers share
one engine and an engine fix reaches all ten.
The scripts run from the workdir and read its README and flags;
nothing is configured twice.}
\label{fig:sys}
\end{figure*}
